% 九、灵敏度分析与组件比较
\section{灵敏度分析与组件比较}\label{sec:sensitivity}
本章回答两个问题：主方法对超参数是否敏感（\S\ref{sec:sens-param}）；各项组件变化怎样影响费用（\S\ref{sec:ablation}）。全部实验与主结果同一 walk-forward 口径、同一执行器、连续执行全年。

\subsection{参数灵敏度}\label{sec:sens-param}
\textbf{场景数 $M$}。固定其余配置（48h+$V(e)$ 主管线），$M\in\{6,12,20,30\}$ 的全年费用为 1375.07、1358.54、1350.75、1352.96 万元：从 6 到 20 条单调改善约 24 万元（更多场景更完整地刻画残差分布，保留水平不再被少数极端路径左右），20 之后进入 $\pm3$ 万元的平台期。$M=30$ 与 $M=20$ 差 0.16\%，说明主结果不依赖场景数的精细选择；取 $M=30$ 兼顾平台稳定与求解时间。

\textbf{前瞻时域}。$M=12$ 下 24h 线性续存与 48h+$V(e)$ 的全年费用为 1362.07 与 1358.54 万元，48h 组件净省 3.53 万元；该增益在 $M=30$ 下为 3.68 万元（表~\ref{tab:components}），两个场景数下方向与量级一致。机理：线性续存费率把午夜库存按固定单价 $v$ 折现，而 $V(e)$ 是凸函数——库存越高边际价值越低；日末光伏富余大的日子，线性口径会高估囤电价值、诱导多买，$V(e)$ 修正了这一偏差（24h 与 48h 配置的平均 24:00 库存分别为 9245.64\,kWh、7200.17\,kWh，不同末端定价对应不同的库存结转策略）。

\textbf{历史窗参数}。续存费率 $v=0.481548$ 并非自由参数，而是“次日 0--5 时均价 $0.4334/\eta$”的推导量；组合权重 $w=0.388815$ 由一月配对误差按最小二乘因果估计（式\eqref{eq:q3combine}），二者都有明确来源且不使用 2--12 月信息。消融配置的 $K_L=6$、$K_P=7$ 与电价窗 $K_p=35$ 由一月实际费用网格搜索选定——全部参数选择只使用一月数据，评分期（2--12 月）从未参与任何调参，杜绝信息泄漏。

\begin{figure}[!htbp]\centering
\safeincludegraphics[width=0.96\textwidth]{fig_f13_sensitivity}
\caption{问题二场景数与前瞻时域的费用比较}\label{fig:fig_f13_sensitivity}
\end{figure}

\subsection{组件比较}\label{sec:ablation}
各项变体均与对应分支的 24h 基线比较，费用增量为“变体费用减基线费用”，见图~\ref{fig:fig_f18_components} 与表~\ref{tab:components}。其中既有启用组件，也有替换预测或场景方法，不把这些实验描述为从 48h 主方法逐项移除组件。负值表示变体更省费，正值表示变体更贵。逐组件解读：

\textbf{48h+$V(e)$（四分支 $-2.4$ 至 $-4.1$ 万元）}：在全部分支都带来节省的时域组件，理由见上文时域段；作为整体组件进入主方法。

\textbf{分解式负载预测（对照“旧均值”，四分支 $+14.5$ 至 $+21.2$ 万元）}：贡献最大的单一组件。日类型均值预测抹平了“周五、周六低负载”的水平切换，而分解式预测用式\eqref{eq:q2beta} 的中位数比率显式建模切换，预测残差方差更小、场景带更窄、保留水平更准。

\textbf{全历史残差池（对照“同类型池”，四分支 $+4.2$ 至 $+5.8$ 万元）}：只用同类型日残差会使早期可用路径减半，场景数不足的代价超过了类型匹配带来的分布贴合收益；残差（预测误差）本身的类型依赖性远弱于负载水平的类型依赖性——切换信息已被点预测吸收。

\textbf{候选计划重评（问 2 增量 $0.00$、问 3 $-0.20$ 万元，问 4-3 为 $-34.96$ 元）}：在 $M=30$ 的这组 24h 对照中，启用候选重评带来的费用变化很小，问 2 与问 4-2 的账单差为零；其收益不能与 48h 组件简单相加，也不计为当前 48h 主结果的费用来源。

\textbf{新价格预测（问 4 对照“旧价格”，4-2 $+0.55$、4-3 $+0.38$ 万元）}：增益存在但有限，与问题四的观察一致——形状易预测导致简单水平规则已能拿到大部分价值。

以上为本评分期上的回顾性组件比较：它支持在该分支保留相应组件，不构成组件收益跨年份可加或普适的证明。


\begin{figure}[H]\centering
\safeincludegraphics[width=0.90\textwidth]{fig_f18_components}
\caption{组件变化的费用增量（相对各分支 24h 基线）}\label{fig:fig_f18_components}
\end{figure}

\begin{table}[H]\centering
\caption{组件变化相对对应单日时域配置的费用增量（元）}\label{tab:components}
\small\renewcommand{\arraystretch}{1.12}\setlength{\tabcolsep}{4pt}
\begin{tabularx}{\textwidth}{@{}>{\hsize=1.5\hsize\linewidth=\hsize}L>{\hsize=.875\hsize\linewidth=\hsize}R>{\hsize=.875\hsize\linewidth=\hsize}R>{\hsize=.875\hsize\linewidth=\hsize}R>{\hsize=.875\hsize\linewidth=\hsize}R@{}}
\toprule
\rowcolor{cumcmgray} 组件变化 & \shortstack{问 2\\单日 13566396.26} & \shortstack{问 3\\单日 13024596.85} & \shortstack{问 4-2\\单日 14270606.71} & \shortstack{问 4-3\\单日 13716972.04} \\
\midrule
启用 48h 与次日价值 & -36761.01 & -35046.21 & -40706.98 & -23960.45 \\
启用候选计划重评 & 0.00 & -1974.51 & 0.00 & -34.96 \\
改用同类型日场景池 & 42116.91 & 47354.21 & 57738.37 & 53857.16 \\
改用旧均值负载预测 & 204242.11 & 145252.73 & 211566.29 & 144922.98 \\
改用旧价格预测 & --- & --- & 5524.96 & 3816.14 \\
\bottomrule
\end{tabularx}
\par\vspace{3pt}{\footnotesize 费用增量为变体减单日时域配置，负值表示节省；”---”表示该分支不适用。}
\end{table}

